\chapter{L'économie d'une société de trading}\label{ch:fm:the-economics-of-a-trading-firm}

En 2020, année de forte volatilité, le chiffre d'affaires d'un teneur de marché électronique coté est passé de 1\,517 millions de dollars à 3\,239 millions de dollars, soit 2.13 fois le
niveau de l'année précédente. Sa rémunération a progressé de 2.6~\% et sa facture de communication et de données de 2.1~\%. La même base de coûts, qui l'avait laissée avec une perte avant impôt de 116 millions de dollars
en 2019, lui a laissé un bénéfice avant impôt de 1\,383 millions de dollars en 2020. Une société de gestion quantitative cotée, rémunérée sur les actifs qu'elle gère et sur la
performance qu'elle délivre, a connu un cycle différent sur les mêmes années~: son chiffre d'affaires montait quand ses fonds de suivi de tendance (trend following) performaient et baissait quand ils
ne performaient pas, quelle que soit l'évolution de l'indice de volatilité, et elle reversait en primes une large part de chaque dollar supplémentaire. Ce chapitre examine ces deux entreprises ainsi que la
division marchés d'une banque, à travers un même jeu de lignes standard, et pose à chacune les mêmes deux questions~: quels coûts sont fixes, et de combien
le chiffre d'affaires peut baisser avant que le bénéfice ne disparaisse.

\section{Le compte de résultat d'une société de trading}

Toute société de trading, quelle que soit la façon dont elle nomme ses lignes, répartit son chiffre d'affaires entre les mêmes quatre postes~: des coûts qui croissent avec l'activité qu'elle réalise (frais de bourse
et de compensation, paiements pour flux d'ordres, financement des positions qu'elle détient), sur la rémunération, sur l'infrastructure fixe de l'entreprise (technologie,
données de marché, locaux, administration), et sur ce qui revient à ses propriétaires et à l'administration fiscale.

\begin{definition}[Revenu net de trading, coûts liés au volume]\label{def:fm:the-economics-of-a-trading-firm:ntr}
Les \emph{coûts liés au volume}\index{coûts liés au volume} d'une société de trading sont les coûts qui croissent proportionnellement à l'activité qu'elle réalise~: frais de bourse,
de compensation et de courtage, paiements pour flux d'ordres, et financement de ses positions et des marges qui les sous-tendent. Son \emph{revenu net
de trading}\index{revenu net de trading} est son revenu de trading et de commissions, y compris les intérêts et dividendes perçus sur les positions, diminué de ses
coûts liés au volume.
\end{definition}

Le \omterm{def:fm:the-economics-of-a-trading-firm:ntr}{revenu net de trading} est le taux de capture du revenu (revenue capture) du teneur de marché (livre~11, chapitre~1)
cumulé sur le volume d'une année, et le bon point de départ pour comparer les entreprises~: deux entreprises qui gagnent la même somme après frais font le même métier, quel que
soit leur chiffre d'affaires brut. Un teneur de marché qui détient de gros inventaires financés par du prêt de titres déclare à la fois les intérêts qu'il perçoit et les
intérêts qu'il verse~; les compenser est ce qui rend son revenu comparable aux frais perçus par une société de gestion.

\begin{definition}[Base de coûts fixes]\label{def:fm:the-economics-of-a-trading-firm:fixed}
La \emph{base de coûts fixes}\index{base de coûts fixes} d'une entreprise est la part de ses coûts qui ne varie pas avec son chiffre d'affaires sur une année~: la rémunération fixe,
technologie et communication, données de marché, locaux, administration et amortissements. Sa \emph{marge opérationnelle}\index{marge opérationnelle} est
le bénéfice opérationnel, c'est-à-dire le \omterm{def:fm:the-economics-of-a-trading-firm:ntr}{revenu net de trading} diminué de la rémunération et des autres coûts fixes, rapporté au \omterm{def:fm:the-economics-of-a-trading-firm:ntr}{revenu net de trading}.
\end{definition}

\begin{omfigure}
\begin{tikzpicture}[font=\footnotesize,>=stealth,
  box/.style={draw,rounded corners,minimum width=3.2cm,minimum height=0.62cm,align=center},
  cost/.style={box,fill=omThm!10},total/.style={box,fill=omDef!12}]
\node[total] (rev) at (0,4.2) {revenu};
\node[cost] (vol) at (5.2,4.2) {coûts liés au volume\\{\scriptsize frais, compensation,}\\{\scriptsize flux d'ordres, financement}};
\node[total] (ntr) at (0,2.8) {revenu net de trading};
\node[cost] (pay) at (5.2,2.8) {rémunération\\{\scriptsize partie fixe + partie variable}};
\node[cost] (fix) at (5.2,1.4) {autres coûts fixes\\{\scriptsize technologie, données,}\\{\scriptsize locaux, administration}};
\node[total] (op) at (0,1.4) {bénéfice opérationnel};
\node[box,text width=4.0cm,align=center] (below) at (0,0) {éléments exceptionnels,\\intérêts de la dette, impôt};
\node[total] (ni) at (5.2,0) {résultat net};
\draw[->] (rev) -- (ntr); \draw[->] (ntr) -- (op); \draw[->] (op) -- (below); \draw[->] (below) -- (ni);
\draw[->,omThm] (vol.west) -- ++(-1.1,0) |- (ntr.east);
\draw[->,omThm] (pay.west) -- ++(-1.1,0) |- (op.north east);
\draw[->,omThm] (fix.west) -- (op.east);
\node[font=\scriptsize,omThm,anchor=west] at (7.9,4.2) {varie avec le volume};
\node[font=\scriptsize,omThm,anchor=west] at (7.9,2.8) {varie en partie avec le revenu};
\node[font=\scriptsize,omThm,anchor=west] at (7.9,1.4) {fixe sur une année};
\end{tikzpicture}

\omcaption{Les lignes standard de ce livre. Les lignes publiées par chaque déclarant sont reprises dans celles-ci (\texttt{firm.firmecon})~; les coûts de droite sont
retranchés dans l'ordre, et l'effet de chacun quand le chiffre d'affaires varie est indiqué à côté.}\label{fig:fm:the-economics-of-a-trading-firm:lines}
\end{omfigure}

La correspondance est un tableau qui relie les lignes publiées par chaque déclarant aux lignes standard (voir le \cref{lst:fm:the-economics-of-a-trading-firm:map}). Trois déclarants de
trois modèles d'activité traversent le chapitre~: les rapports annuels d'un teneur de marché électronique, déposés sur le Form 10-K pour 2019--2025, et ceux d'une société de gestion quantitative
cotée, pour 2020--2025, et le segment marchés d'une banque pour 2023--2025. Leurs chiffres, en millions de dollars tels que déposés, sont consignés dans de petits tableaux dérivés
avec le code du chapitre.

\omcode{code/desk/01-the-economics-of-a-trading-firm/python/fm_econ.py}{23}{40}{Les tableaux de correspondance~: chaque ligne standard est une somme algébrique de
lignes publiées par un déclarant.}\label{lst:fm:the-economics-of-a-trading-firm:map}

\begin{example}[Le revenu net de trading du teneur de marché en 2025]\label{ex:fm:the-economics-of-a-trading-firm:virtu}
Le compte de résultat 2025 de Virtu Financial fait apparaître un revenu de trading de 2\,436.7 millions de dollars, des commissions et services technologiques de 617.0 millions de dollars,
des produits d'intérêts et de dividendes de 508.8 millions de dollars et, parmi ses charges, les frais de courtage, de bourse et de compensation ainsi que les paiements pour flux d'ordres de
769.8 millions de dollars, et des charges d'intérêts et de dividendes de 647.4 millions de dollars. Son \omterm{def:fm:the-economics-of-a-trading-firm:ntr}{revenu net de trading} est $2\,436.7+617.0+508.8-769.8-647.4=\$2\,145.3$
millions de dollars~; les \omterm{def:fm:the-economics-of-a-trading-firm:ntr}{coûts liés au volume} absorbent $1\,417.2/3\,562.5=39.8\%$ du revenu avant autres produits. La rémunération (528.1 millions de dollars) et les autres
coûts fixes (communication et données 249.2 millions de dollars, opérations et administration 97.9 millions de dollars, dotations aux amortissements 64.4 millions de dollars, 411.5
millions de dollars au total) laissent un bénéfice opérationnel de 1\,205.7 millions de dollars, soit une marge de 56.2~\% sur le \omterm{def:fm:the-economics-of-a-trading-firm:ntr}{revenu net de trading}.
\end{example}

L'entreprise publie elle-même le même montant, 2\,145.3 millions de dollars, sous un autre nom, le \omterm{def:fm:the-economics-of-a-trading-firm:ntr}{revenu net de trading} ajusté, et le qualifie d'\omterm{def:fm:the-economics-of-a-trading-firm:nongaap}{indicateur non GAAP}. La plupart des sociétés de trading cotées
publient de tels indicateurs, parce que les normes comptables n'ont pas été écrites pour leur activité~; la tâche du lecteur est de retrouver la
réconciliation et vérifier que les ajustements jouent dans les deux sens.

\begin{definition}[Indicateur non GAAP]\label{def:fm:the-economics-of-a-trading-firm:nongaap}
Un \emph{indicateur non GAAP}\index{indicateur non GAAP} est un chiffre de performance ou de situation qu'une entreprise publie à côté de ses
comptes légaux, mais qui exclut des montants que les normes comptables incluent, ou qui inclut des montants qu'elles excluent~: chiffre d'affaires ajusté, bénéfice courant, résultat
ajusté. Les régulateurs des marchés financiers exigent que cette mesure soit étiquetée, définie et réconciliée avec le chiffre légal le plus proche.
\end{definition}

\begin{dated}{2026-09}{Les règles sur les indicateurs non GAAP et alternatifs}\label{dat:fm:the-economics-of-a-trading-firm:rules}
Aux États-Unis, la Regulation G (17 CFR 244.100--102) exige d'un émetteur qui publie un indicateur financier non GAAP qu'il présente la mesure GAAP la plus
directement comparable avec elle et une réconciliation quantitative des deux~; l'Item 10(e) de la Regulation S-K s'applique de la même façon aux dépôts auprès
de la SEC. Dans l'Union européenne, les Guidelines on Alternative Performance Measures de l'ESMA (ESMA/2015/1415, 5 octobre 2015) exigent que ces mesures soient
définies, dotées de libellés explicites et réconciliées avec la ligne la plus directement rapprochable des états financiers. Man Group, par exemple, indique
que ses mesures «~core~» sont des mesures alternatives de performance et les réconcilie avec le résultat légal.
\end{dated}

\section{Coûts qui croissent et coûts qui ne croissent pas}

Les trois entreprises répartissent leur chiffre d'affaires selon des proportions très différentes (voir la \cref{fig:fm:the-economics-of-a-trading-firm:shares}). Le teneur de marché verse
quatre dollars sur dix aux bourses, aux chambres de compensation, aux courtiers et aux prêteurs avant de payer quiconque~; les \omterm{def:fm:the-economics-of-a-trading-firm:ntr}{coûts liés au volume} de la société de gestion, son
administration des actifs, sont de cinq pour cent, mais elle verse presque la moitié de son chiffre d'affaires à son personnel~; le segment marchés de la banque ne déclare aucune ligne liée au volume qui lui soit propre,
car ses commissions figurent dans les autres charges d'exploitation, et celles-ci sont plus élevées que sa rémunération.

\begin{omfigure}
\begin{tikzpicture}
\begin{axis}[width=11cm,height=5.4cm,xbar stacked,bar width=13pt,xmin=0,xmax=100,xlabel={\small part du revenu (\%)},
  ytick={0,1,2},yticklabels={teneur de marché,société de gestion,marchés de banque},y dir=reverse,ymin=-0.6,ymax=2.6,
  tick label style={font=\footnotesize},table/col sep=comma,grid=major,grid style={gray!20},
  legend style={font=\scriptsize,at={(0.5,-0.32)},anchor=north,/tikz/every even column/.append style={column sep=8pt}},legend columns=4]
\addplot[fill=omThm!55,draw=omThm] table[y=k,x=volume]{figdata/desk/01-the-economics-of-a-trading-firm/shares.csv};
\addplot[fill=omProp!55,draw=omProp] table[y=k,x=comp]{figdata/desk/01-the-economics-of-a-trading-firm/shares.csv};
\addplot[fill=gray!35,draw=gray] table[y=k,x=fixed]{figdata/desk/01-the-economics-of-a-trading-firm/shares.csv};
\addplot[fill=omDef!55,draw=omDef] table[y=k,x=profit]{figdata/desk/01-the-economics-of-a-trading-firm/shares.csv};
\legend{coûts liés au volume,rémunération,autres coûts fixes,bénéfice opérationnel}
\end{axis}
\end{tikzpicture}

\omcaption{Répartition du revenu 2025 de chaque entreprise, en pourcentage du revenu avant autres produits~: un teneur de marché électronique (Virtu Financial), une société de gestion cotée
quantitative (Man Group, mesures centrales) et le segment marchés d'une banque (Goldman Sachs, Global Banking \& Markets, net des
provisions pour créances douteuses). Données~: Forms 10-K 2025 et résultats 2025 de Man Group, via \texttt{fm\_econ.latest\_shares}.}\label{fig:fm:the-economics-of-a-trading-firm:shares}
\end{omfigure}

La répartition entre les coûts qui varient avec l'activité et ceux qui ne varient pas détermine ce que produit une mauvaise année pour une entreprise. Notons le \omterm{def:fm:the-economics-of-a-trading-firm:ntr}{revenu net de trading} $\mathrm{Rev}$ (net des
\omterm{def:fm:the-economics-of-a-trading-firm:ntr}{coûts liés au volume}), la part fixe des coûts, rémunération fixe comprise, $C^{\mathrm{fix}}$, et posons que la rémunération variable prend une part $b$ de chaque
dollar de revenu net. Le bénéfice opérationnel vaut
\[
\Pi(\mathrm{Rev})=(1-b)\,\mathrm{Rev}-C^{\mathrm{fix}} .
\]

\begin{proposition}[Baisse du seuil de rentabilité et levier opérationnel]\label{prop:fm:the-economics-of-a-trading-firm:fall}
Si $\Pi(\mathrm{Rev})>0$, le bénéfice opérationnel s'annule lorsque le \omterm{def:fm:the-economics-of-a-trading-firm:ntr}{revenu net de trading} baisse de la fraction
\[
x^\ast=1-\frac{C^{\mathrm{fix}}}{(1-b)\,\mathrm{Rev}},
\]
et l'élasticité du bénéfice opérationnel au \omterm{def:fm:the-economics-of-a-trading-firm:ntr}{revenu net de trading}, c'est-à-dire le levier opérationnel, vaut
\[
\frac{d\ln\Pi}{d\ln\mathrm{Rev}}=\frac{(1-b)\,\mathrm{Rev}}{(1-b)\,\mathrm{Rev}-C^{\mathrm{fix}}}=\frac1{x^\ast}.
\]
Pour des coûts totaux donnés au revenu courant, une part variable $b$ plus élevée relève la baisse du seuil de rentabilité et abaisse le levier opérationnel.
\end{proposition}

\begin{proof}
Résolvons $(1-b)(1-x)\mathrm{Rev}=C^{\mathrm{fix}}$ en $x$~; dérivons $\ln\Pi$. Pour la comparaison, fixons les coûts totaux
$K=C^{\mathrm{fix}}+b\,\mathrm{Rev}$~: alors $\Pi=\mathrm{Rev}-K$ ne dépend pas de $b$, tandis que $(1-b)\mathrm{Rev}$, le numérateur du
levier, diminue quand $b$ augmente, et $x^\ast=\Pi/((1-b)\mathrm{Rev})$ augmente.
\end{proof}

Le levier opérationnel est la proposition du livre~11 sur le volume d'un teneur de marché (chapitre~1 de ce livre), transposée à l'ensemble de l'entreprise~: là, les coûts fixes portent
sur le volume, ici ils portent sur le revenu net après tous les \omterm{def:fm:the-economics-of-a-trading-firm:ntr}{coûts liés au volume}. La seule nouveauté tient à la rémunération. Une entreprise qui verse une part fixe de chaque
dollar marginal sous forme de primes partage ses mauvaises années avec son personnel~; une entreprise dont la rémunération est surtout constituée de salaires ne le fait pas.

\begin{example}[Une entreprise de 100]\label{ex:fm:the-economics-of-a-trading-firm:small}
Une entreprise dispose d'un \omterm{def:fm:the-economics-of-a-trading-firm:ntr}{revenu net de trading} de 100, d'une rémunération fixe de 15, d'une rémunération variable égale à 30~\% du revenu net et d'autres coûts fixes de 35~: son bénéfice opérationnel vaut
$100-15-30-35=20$. Avec une rémunération qui s'ajuste, $C^{\mathrm{fix}}=50$ et $b=0.3$~: la baisse du seuil de rentabilité vaut $1-50/70=28.6~\%$ et le levier opérationnel 3.5.
Si les mêmes 45 de rémunération étaient tous en salaires, $C^{\mathrm{fix}}=15+30+35=80$ et $b=0$~: la baisse du seuil de rentabilité vaut $1-80/100=20~\%$ et le
levier opérationnel 5. Une rémunération flexible transforme une entreprise qui
perd de l'argent après une baisse de 20~\% en une entreprise qui survit à une baisse de 28~\%.
\end{example}

\begin{method}[Lire la structure de coûts d'une entreprise à partir de ses dépôts réglementaires]\label{met:fm:the-economics-of-a-trading-firm:read}
\begin{enumerate}
\item Reporter chaque ligne publiée sur les lignes standard et vérifier qu'elles font bien le total du résultat avant impôt publié, hors éléments exceptionnels.
\item Calculer le \omterm{def:fm:the-economics-of-a-trading-firm:ntr}{revenu net de trading}, la rémunération et les autres coûts fixes pour chaque année disponible (cinq au minimum).
\item Régresser la rémunération sur le \omterm{def:fm:the-economics-of-a-trading-firm:ntr}{revenu net de trading} au fil des années~: la pente estime $b$, l'ordonnée à l'origine la rémunération fixe. Lire l'erreur type de la pente avant
la pente.
\item Prendre les autres coûts fixes à leur dernier niveau, puis calculer $x^\ast$ et le levier opérationnel au revenu de la dernière année, avec le
$b$ estimé et avec $b=0$.
\end{enumerate}
\end{method}

\section{Rémunération~: le poste le plus lourd, et la part qui est variable}

\begin{definition}[Ratio de rémunération]\label{def:fm:the-economics-of-a-trading-firm:comp}
Le \emph{ratio de rémunération}\index{ratio de rémunération} d'une entreprise est sa rémunération totale, salaires, primes, attributions différées et charges sociales sur la masse salariale, rapportée
à son revenu net sur la même période. Les entreprises qui le publient précisent leur propre définition du revenu net.
\end{definition}

Un \omterm{def:fm:the-economics-of-a-trading-firm:comp}{ratio de rémunération} n'est pas une structure de coûts. Le ratio de la société de gestion a été de 40~\% en 2021 et en 2022, et de 50~\% en 2023 (chiffres propres à Man Group,
rémunération rapportée au revenu net central avant coûts de services aux actifs)~; sa politique affichée est un ratio «~compris entre 40~\% et 50~\% du revenu net central, selon la
composition et le niveau du revenu~». Le ratio a augmenté en 2023 parce que le revenu a baissé plus que la rémunération~: un ratio évolue à l'inverse du revenu dès qu'une partie de la rémunération est
fixe. La pente de la rémunération par rapport au revenu, d'une année à l'autre, indique quelle part de la rémunération est en réalité variable
(la \cref{fig:fm:the-economics-of-a-trading-firm:pay}).

\begin{omfigure}
\begin{tikzpicture}
\begin{axis}[width=10.5cm,height=5.8cm,xlabel={\small revenu net de trading (millions de \$)},ylabel={\small rémunération (millions de \$)},
  xmin=800,xmax=2400,ymin=330,ymax=720,tick label style={font=\footnotesize},scaled ticks=false,
  xtick={800,1200,1600,2000,2400},xticklabels={800,1200,1600,2000,2400},
  grid=major,grid style={gray!20},table/col sep=comma,legend style={font=\scriptsize,at={(0.5,-0.3)},anchor=north,/tikz/every even column/.append style={column sep=8pt}},legend columns=2]
\addplot[only marks,mark=*,omThm] table[x=nr,y=comp]{figdata/desk/01-the-economics-of-a-trading-firm/pay_virtu.csv};
\addplot[omThm,thick] table[x=nr,y=fit]{figdata/desk/01-the-economics-of-a-trading-firm/pay_virtu.csv};
\addplot[only marks,mark=square*,omDef] table[x=nr,y=comp]{figdata/desk/01-the-economics-of-a-trading-firm/pay_man.csv};
\addplot[omDef,thick,dashed] table[x=nr,y=fit]{figdata/desk/01-the-economics-of-a-trading-firm/pay_man.csv};
\legend{teneur de marché (2019--2025),sa droite d'ajustement~: pente 0.05,société de gestion (2020--2025),sa droite d'ajustement~: pente 0.29}
\end{axis}
\end{tikzpicture}

\omcaption{Rémunération en fonction du \omterm{def:fm:the-economics-of-a-trading-firm:ntr}{revenu net de trading}, un point par année, avec la droite des moindres carrés de chaque entreprise. La rémunération du teneur de marché varie à peine
avec le revenu (pente 0.048, erreur type 0.045)~; celle de la société de gestion augmente de 29 cents par dollar (erreur type 0.094). Données~: Virtu Financial
Forms 10-K 2019--2025, résultats de Man Group 2020--2025, via \texttt{fm\_econ.pay\_fit}.}\label{fig:fm:the-economics-of-a-trading-firm:pay}
\end{omfigure}

La rémunération du teneur de marché s'est située entre 376 et 394 millions de dollars chaque année de 2019 à 2023, tandis que son \omterm{def:fm:the-economics-of-a-trading-firm:ntr}{revenu net de trading} a varié de
975 millions à 2\,271 millions de dollars~; la pente vaut 0.048 avec une erreur type de 0.045, ce qui ne se distingue pas de zéro. Sa rémunération 2025, 528 millions de dollars, se place
loin au-dessus de la droite, premier signe qu'elle avait commencé à changer. La pente de la société de gestion, 0.29 avec une erreur type de 0.094, indique
qu'environ trois dixièmes de chaque dollar marginal de son revenu sont allés à la rémunération. Les trois années du segment bancaire donnent une pente de 0.20~; trois points font une droite, mais
pas une estimation. La société de gestion publie aussi directement la répartition pour 2021~: rémunération fixe de 208 millions de dollars, rémunération variable de 388 millions de dollars.

\begin{remark}[Ce que six points permettent de dire]\label{rem:fm:the-economics-of-a-trading-firm:six}
Une pente estimée à partir de six ou sept points annuels ne dispose que de quatre ou cinq degrés de liberté et varie beaucoup dès qu'une année est retirée~: sans 2020, la
pente de la société de gestion passe de 0.29 à 0.13, avec une erreur type de 0.13 (exercice~7). La croissance des effectifs, une acquisition ou un changement dans la
composition des commissions modifie la rémunération pour des raisons étrangères au revenu de l'année. La régression est une première lecture de la structure de coûts, à
confronter à ce que l'entreprise dit de sa rémunération fixe et variable, jamais une mesure de celle-ci.
\end{remark}

\begin{center}
\small
\begin{tabular}{@{}lrrr@{}}
\toprule
2025 & teneur de marché & société de gestion & marchés de banque\\
\midrule
\omterm{def:fm:the-economics-of-a-trading-firm:ntr}{revenu net de trading} (millions de \$) & 2\,145.3 & 1\,325 & 41\,075\\
rémunération / \omterm{def:fm:the-economics-of-a-trading-firm:ntr}{revenu net de trading} & 24.6~\% & 50.9~\% & 26.8~\%\\
\omterm{def:fm:the-economics-of-a-trading-firm:fixed}{marge opérationnelle} & 56.2~\% & 32.3~\% & 42.8~\%\\
\omterm{def:fm:the-economics-of-a-trading-firm:ntr}{revenu net de trading} par tête (millions de \$) & 2.09 & 0.77 & --\\
rémunération par tête (millions de \$) & 0.51 & 0.39 & --\\
\bottomrule
\end{tabular}
\end{center}

\section{Capital, financement et rentabilité des fonds propres}

\begin{definition}[Rentabilité des fonds propres]\label{def:fm:the-economics-of-a-trading-firm:roe}
La \emph{rentabilité des fonds propres}\index{rentabilité des fonds propres} (return on equity), ROE, d'une entreprise (ou d'une activité) est son résultat net sur une période divisé par la moyenne de ses
capitaux propres sur la même période~; pour une division d'un groupe plus vaste, les fonds propres sont la part des fonds propres du groupe que l'entreprise lui attribue.
\end{definition}

Le teneur de marché a dégagé un résultat net, intérêts minoritaires compris, de 912.3~millions de dollars en 2025 sur des capitaux propres totaux moyens de 1\,730~millions de dollars~:
une \omterm{def:fm:the-economics-of-a-trading-firm:roe}{rentabilité des fonds propres} de 52.7~\%. En 2023, année calme, le même calcul donne 17.3~\%. Le segment banque publie une rentabilité moyenne des capitaux propres ordinaires de 16.4~\%
pour 2025, sur des fonds propres alloués de 79.7~milliards de dollars. Les deux chiffres ne mesurent pas la même chose. Les fonds propres d'un teneur de marché sont ce que ses propriétaires ont
laissé dans l'entreprise~; elle est soumise aux règles de capital applicables aux courtiers et aux entreprises d'investissement (chapitre~2), et le bilan qu'elle finance avec l'argent de ses
courtiers et de ses prêteurs est bien plus grand que ses fonds propres. Les fonds propres de segment d'une banque sont une allocation, fixée par celle des contraintes de capital du groupe
qui est saturée (chapitre~5), et le groupe en détient davantage que son activité de marché ne pourrait perdre en un an, par construction réglementaire.

\begin{dated}{2026-09}{Trois entreprises dans leurs dépôts de 2025}\label{dat:fm:the-economics-of-a-trading-firm:three}
\textbf{Virtu Financial} (formulaire 10-K pour 2025)~: produit total de 3\,632.1~millions de dollars~; lignes liées au volume de 1\,417.2~millions de dollars (courtage, frais de bourse,
frais de compensation et paiements pour flux d'ordres de 769.8~millions de dollars, charges d'intérêts et de dividendes de 647.4~millions de dollars)~; rémunération des salariés de 528.1~millions de dollars~;
\omterm{def:fm:the-economics-of-a-trading-firm:ntr}{revenu net de trading} ajusté, \omterm{def:fm:the-economics-of-a-trading-firm:nongaap}{indicateur non GAAP}, de 2\,145.3~millions de dollars~; environ 1\,027 salariés en février 2026.
\textbf{Man Group} (résultats pour 2025)~: produit net cœur de métier de 1\,398~millions de dollars, dont frais de gestion nets de 1\,077~millions de dollars et commissions de performance de 281
millions de dollars~; rémunération cœur de métier de 675~millions de dollars, soit un \omterm{def:fm:the-economics-of-a-trading-firm:comp}{ratio de rémunération} de 48~\% (47~\% en 2024)~; résultat avant impôt cœur de métier de 407~millions de dollars~; 1\,719 collaborateurs à la fin
de l'année. \textbf{Goldman Sachs, Global Banking \& Markets} (formulaire 10-K pour 2025)~: produits nets de 41\,453~millions de dollars~; rémunérations et avantages du personnel
de 11\,025~millions de dollars~; autres charges d'exploitation de 12\,476~millions de dollars~; résultat avant impôt de 17\,574~millions de dollars~; capitaux propres ordinaires moyens de 79\,748~millions de dollars~; rentabilité
moyenne des capitaux propres ordinaires de 16.4~\% (13.8~\% en 2024, 11.3~\% en 2023).
\end{dated}

\begin{remark}[Le financement est un coût lié au volume]\label{rem:fm:the-economics-of-a-trading-firm:financing}
Entre 2021 et 2025, les charges d'intérêts et de dividendes du teneur de marché sont passées de 140~millions à 647~millions de dollars, et ses produits
d'intérêts et de dividendes de 75~millions à 509~millions de dollars~: des taux plus élevés ont relevé les deux côtés du financement de son inventaire. Net, la variation reste faible~; brute, elle
a plus que doublé les \omterm{def:fm:the-economics-of-a-trading-firm:ntr}{coûts liés au volume} publiés. C'est pourquoi ce livre intègre le financement dans le \omterm{def:fm:the-economics-of-a-trading-firm:ntr}{revenu net de trading}, et pourquoi le produit publié
d'une entreprise peut augmenter avec les taux d'intérêt sans que son activité ne croisse.
\end{remark}

\section{Le cycle~: volatilité, volume et revenu}

Le revenu net d'un teneur de marché est égal au volume multiplié par le taux de capture (livre~11), et les deux augmentent avec la volatilité. Une régression du \omterm{def:fm:the-economics-of-a-trading-firm:ntr}{revenu net de trading}
du teneur de marché sur la moyenne annuelle de l'indice VIX donne une pente de 52~millions de dollars par point d'indice (erreur type~33) et une corrélation de 0.58 sur sept
ans. Le revenu net de la société de gestion présente une corrélation de $-0.27$ avec le même indice~: son revenu suit ses actifs et la performance de ses fonds,
non la volatilité du marché. Sur les sept années, la rémunération et les coûts fixes du teneur de marché ont évolué dans une bande étroite, tandis que le revenu
a varié d'un facteur 2.3 (la \cref{fig:fm:the-economics-of-a-trading-firm:virtu}).

\begin{omfigure}
\begin{tikzpicture}
\begin{axis}[width=11cm,height=5.4cm,xlabel={\small année},ylabel={\small millions de \$},xmin=2018.6,xmax=2025.4,ymin=0,ymax=2500,
  xtick={2019,2020,2021,2022,2023,2024,2025},xticklabel style={/pgf/number format/1000 sep={}},scaled ticks=false,
  tick label style={font=\footnotesize},grid=major,grid style={gray!20},table/col sep=comma,
  legend style={font=\scriptsize,at={(0.5,-0.3)},anchor=north,/tikz/every even column/.append style={column sep=8pt}},legend columns=4]
\addplot[omDef,thick,mark=*] table[x=year,y=nr]{figdata/desk/01-the-economics-of-a-trading-firm/virtu_years.csv};
\addplot[omProp,thick,mark=square*] table[x=year,y=comp]{figdata/desk/01-the-economics-of-a-trading-firm/virtu_years.csv};
\addplot[gray,thick,mark=triangle*] table[x=year,y=fixed]{figdata/desk/01-the-economics-of-a-trading-firm/virtu_years.csv};
\addplot[omThm,thick,dashed,mark=diamond*] table[x=year,y=profit]{figdata/desk/01-the-economics-of-a-trading-firm/virtu_years.csv};
\legend{revenu net de trading,rémunération,autres coûts fixes,résultat opérationnel}
\end{axis}
\end{tikzpicture}

\omcaption{Le teneur de marché, 2019--2025~: \omterm{def:fm:the-economics-of-a-trading-firm:ntr}{revenu net de trading} (revenu diminué des frais, des paiements pour flux d'ordres et du financement), rémunération, autres coûts fixes et
résultat opérationnel (avant amortissement des actifs incorporels, intérêts de la dette, éléments exceptionnels et impôt). Données~: formulaires 10-K de Virtu Financial pour 2019--2025, via
\texttt{fm\_econ.table}.}\label{fig:fm:the-economics-of-a-trading-firm:virtu}
\end{omfigure}

En appliquant la \cref{met:fm:the-economics-of-a-trading-firm:read} à la dernière année de chaque entreprise, on obtient sa baisse au seuil de rentabilité et son levier opérationnel,
la rémunération variant selon la pente estimée et la rémunération étant maintenue à son dernier niveau~:

\begin{center}
\small
\begin{tabular}{@{}lrrrr@{}}
\toprule
 & \multicolumn{2}{c}{rémunération variable} & \multicolumn{2}{c}{rémunération fixe}\\
 & baisse au seuil & levier & baisse au seuil & levier\\
\midrule
teneur de marché (2025) & 59.0~\% & 1.69 & 56.2~\% & 1.78\\
société de gestion (2025) & 45.6~\% & 2.19 & 32.3~\% & 3.10\\
marchés de banque (2025) & 53.3~\% & 1.88 & 42.8~\% & 2.34\\
\bottomrule
\end{tabular}
\end{center}

\begin{omfigure}
\begin{tikzpicture}
\begin{axis}[width=10.5cm,height=5.6cm,xlabel={\small variation du revenu net de trading (\%)},ylabel={\small résultat opérationnel (dernière année $=1$)},
  xmin=-60,xmax=40,ymin=-0.4,ymax=2,tick label style={font=\footnotesize},grid=major,grid style={gray!20},table/col sep=comma,
  legend style={font=\scriptsize,at={(0.5,-0.3)},anchor=north,/tikz/every even column/.append style={column sep=8pt}},legend columns=3]
\addplot[omThm,thick] table[x=change,y=virtu]{figdata/desk/01-the-economics-of-a-trading-firm/cycle.csv};
\addplot[omDef,thick,dashed] table[x=change,y=man]{figdata/desk/01-the-economics-of-a-trading-firm/cycle.csv};
\addplot[omProp,thick,dotted] table[x=change,y=gsgbm]{figdata/desk/01-the-economics-of-a-trading-firm/cycle.csv};
\draw[gray] (axis cs:-60,0) -- (axis cs:40,0);
\legend{teneur de marché,société de gestion,marchés de banque}
\end{axis}
\end{tikzpicture}

\omcaption{Résultat opérationnel après une variation du \omterm{def:fm:the-economics-of-a-trading-firm:ntr}{revenu net de trading}, par rapport à celui de la dernière année, la rémunération variant selon la pente estimée de chaque entreprise.
Chaque courbe coupe zéro à sa baisse au seuil~: 59~\% pour le teneur de marché, 46~\% pour la société de gestion, 53~\% pour le segment banque. Données~:
\texttt{fm\_econ.profit\_curves} sur les dépôts de 2025.}\label{fig:fm:the-economics-of-a-trading-firm:cycle}
\end{omfigure}

Le teneur de marché, l'entreprise à la rémunération la plus rigide, est le plus robuste des trois à son revenu de 2025, parce que sa marge est la plus large~: une marge élevée
représente une grande distance au seuil de rentabilité, quelle que soit la manière dont elle a été obtenue. La rémunération flexible de la société de gestion est ce qui la maintient à flot lors d'une mauvaise année~: maintenue fixe, sa
rémunération ramènerait le seuil à une baisse de 32~\% au lieu de 46~\%. En 2019, la marge du teneur de marché était de 21.8~\%, et sa baisse au seuil, au revenu de
cette année-là, avec la même pente de rémunération, était de 22.9~\%, contre 59.0~\% en 2025. La robustesse d'une entreprise est une propriété de sa structure de coûts \emph{et} de sa position dans le cycle.

\section{Tutoriel~: trois entreprises à travers un même jeu de lignes}\label{tut:fm:the-economics-of-a-trading-firm:tut}

\begin{tutorial}
\textbf{Objectif.} Lire les dépôts de trois entreprises à travers un même plan comptable et mesurer la structure de coûts de chacune. \textbf{État final~:} le tableau des
baisses au seuil et la \cref{fig:fm:the-economics-of-a-trading-firm:cycle}.
\begin{tutsteps}
\item \textbf{Les données.} \texttt{data/desk/filings\_virtu.csv}, \texttt{filings\_man.csv} et \texttt{filings\_gsgbm.csv} contiennent les lignes de chaque déclarant,
en millions de dollars, le dépôt et la page étant consignés dans \texttt{data/desk/LICENSES.md}~; \texttt{vix\_annual.csv} contient la moyenne annuelle de l'indice,
produite par \texttt{fm\_fetch\_vix.py} (la seule étape qui nécessite le réseau).
\item \textbf{Correspondance.} \texttt{fm\_econ.table(firm)} projette chaque année sur les lignes standard grâce à \texttt{firm.firmecon.statements} et aux tableaux du
\cref{lst:fm:the-economics-of-a-trading-firm:map}. Vérifier que le \omterm{def:fm:the-economics-of-a-trading-firm:ntr}{revenu net de trading} du teneur de marché reproduit son propre chiffre non GAAP,
soit 2\,145.3~millions de dollars en 2025 et 1\,597.7~millions de dollars en 2024.
\item \textbf{Ajuster la rémunération.} \texttt{fm\_econ.pay\_fit(firm)} régresse la rémunération sur le \omterm{def:fm:the-economics-of-a-trading-firm:ntr}{revenu net de trading} (la \cref{fig:fm:the-economics-of-a-trading-firm:pay}).
\item \textbf{Seuil de rentabilité.} \texttt{fm\_econ.structure(firm, flex)} construit la structure de coûts et calcule $x^\ast$ et le levier
(le \cref{lst:fm:the-economics-of-a-trading-firm:cs})~; \texttt{profit\_curves()} trace la \cref{fig:fm:the-economics-of-a-trading-firm:cycle}.
\omcode{code/firm/firmecon/firm_firmecon.py}{100}{126}{La structure de coûts à partir des années propres à une entreprise, et la baisse au seuil et le levier
opérationnel de la \cref{prop:fm:the-economics-of-a-trading-firm:fall}.}\label{lst:fm:the-economics-of-a-trading-firm:cs}
\item \textbf{Le cycle.} \texttt{fm\_econ.vix\_fit(firm)} régresse le \omterm{def:fm:the-economics-of-a-trading-firm:ntr}{revenu net de trading} sur la moyenne annuelle de l'indice.
\end{tutsteps}
\end{tutorial}

\textbf{Pour aller plus loin.} Retirer une année à la fois et observer le déplacement de chaque pente (exercice 7)~; ajouter un quatrième déposant à partir de son rapport annuel, en
n'écrivant qu'un nouveau tableau de correspondance.

\section{Construction~: l'économie de l'entreprise en code}\label{bld:fm:the-economics-of-a-trading-firm:firmecon}

\begin{build}
\bfield{Rôle} Un seul plan de comptes pour chaque modèle économique du livre, afin que les chapitres suivants (le partenariat, la plateforme, le fonds, le
desk de banque, le plan d'entrée) énoncent leur économie dans les mêmes lignes.
\bfield{Interface} \texttt{firm.firmecon}~: \texttt{Statement(year, net\_revenue, comp, other\_fixed, volume\_cost, headcount, equity)} avec
\texttt{operating\_profit}~; \texttt{map\_lines(row, mapping)}, \texttt{statements(rows, mapping)}~; \texttt{ratios(st)}~; \texttt{fit\_line(x, y)}
avec des erreurs types~; \texttt{CostStructure(fixed, var\_share)}, \texttt{cost\_structure(sts, flex\_pay)}~; \texttt{profit},
\texttt{breakeven\_fall}, \texttt{operating\_leverage}, \texttt{cycle}.
\bfield{Règles} Un mapping nomme chaque colonne publiée qu'il utilise, avec son signe~; rien n'est estimé à partir de moins de trois années~; le levier opérationnel est
indéfini au seuil de rentabilité ou en deçà, et lève une erreur à ce niveau.
\bfield{Tests d'acceptation} \texttt{code/firm/firmecon/tests/}~: un mapping reproduit un compte rendu construit à la main~; la droite ajustée est exacte sur des données exactes~;
le profit est nul après la baisse au seuil de rentabilité~; le levier égale l'élasticité numérique~; une rémunération flexible abaisse le levier et élève la baisse au seuil
à coût total égal.
\bfield{Pour aller plus loin} Un analyseur des faits XBRL d'un déposant vers les colonnes du mapping~; un scénario de cycle dans lequel les \omterm{def:fm:the-economics-of-a-trading-firm:ntr}{coûts liés au volume}, le revenu net et la
rémunération variable évoluent tous avec une trajectoire de volatilité.
\end{build}

\begin{omsources}
\item Virtu Financial, Inc., Forms 10-K pour 2019 à 2025 (SEC EDGAR), et les XBRL company facts de la SEC pour l'entreprise.
\item Man Group plc, Results for the years ended 31 December 2021, 2023 and 2025.
\item The Goldman Sachs Group, Inc., Form 10-K pour 2025, résultats par segment.
\item Cboe Global Markets, historique quotidien de l'indice VIX (moyennes annuelles dérivées uniquement).
\item US Code of Federal Regulations, 17 CFR Part 244 (Regulation G)~; ESMA, Guidelines on Alternative Performance Measures, ESMA/2015/1415.
\end{omsources}

\section{Exercices}

\begin{exercise}[$\star$]\label{exo:fm:the-economics-of-a-trading-firm:1}
À partir des lignes 2025 du teneur de marché dans l'\cref{ex:fm:the-economics-of-a-trading-firm:virtu}, calculer son \omterm{def:fm:the-economics-of-a-trading-firm:ntr}{revenu net de trading} et la part de son
revenu avant autres produits absorbée par les \omterm{def:fm:the-economics-of-a-trading-firm:ntr}{coûts liés au volume}.
\end{exercise}

\begin{exercise}[$\star$]\label{exo:fm:the-economics-of-a-trading-firm:2}
La société de gestion a versé 675~millions de dollars sur un revenu net de base de 1\,398~millions de dollars en 2025, et 678~millions de dollars sur 1\,696~millions de dollars en 2022. Calculer son
\omterm{def:fm:the-economics-of-a-trading-firm:comp}{ratio de rémunération} chaque année, selon sa propre définition, et expliquer pourquoi le ratio a augmenté alors que la rémunération a baissé.
\end{exercise}

\begin{exercise}[$\star$]\label{exo:fm:the-economics-of-a-trading-firm:3}
Le résultat net revenant aux actionnaires ordinaires du segment banque s'élevait à 13\,117~millions de dollars en 2025, sur des fonds propres ordinaires moyens de 79\,748~millions de dollars. Vérifier sa \omterm{def:fm:the-economics-of-a-trading-firm:roe}{rentabilité
des fonds propres} publiée.
\end{exercise}

\begin{exercise}[$\star\star$]\label{exo:fm:the-economics-of-a-trading-firm:4}
Pour l'entreprise de l'\cref{ex:fm:the-economics-of-a-trading-firm:small}, calculer le résultat opérationnel après une baisse de 25~\% du \omterm{def:fm:the-economics-of-a-trading-firm:ntr}{revenu net de trading}, avec une rémunération
qui s'ajuste, puis avec toute la rémunération fixe.
\end{exercise}

\begin{exercise}[$\star\star$]\label{exo:fm:the-economics-of-a-trading-firm:5}
De 2019 à 2020, le \omterm{def:fm:the-economics-of-a-trading-firm:ntr}{revenu net de trading} du teneur de marché est passé de 974.7~millions de dollars à 2\,271.4~millions de dollars, et son résultat opérationnel de 212.4
millions de dollars à 1\,502.8~millions de dollars. Calculer les deux variations en pourcentage et leur rapport, et comparer ce rapport à son levier opérationnel en 2019.
\end{exercise}

\begin{exercise}[$\star\star$]\label{exo:fm:the-economics-of-a-trading-firm:6}
Sur la \cref{fig:fm:the-economics-of-a-trading-firm:pay}, lequel des points du teneur de marché se situe le plus au-dessus de sa droite, et quelles deux lectures de ce point
faudrait-il vérifier dans le dépôt~?
\end{exercise}

\begin{exercise}[$\star\star\star$]\label{exo:fm:the-economics-of-a-trading-firm:7}
\emph{Programmation.} Réajuster la rémunération de la société de gestion sur son \omterm{def:fm:the-economics-of-a-trading-firm:ntr}{revenu net de trading} sans 2020. Quelles sont la pente et son erreur type, et que dit le
changement sur la fiabilité de l'estimation~?
\end{exercise}

\begin{exercise}[$\star\star\star$]\label{exo:fm:the-economics-of-a-trading-firm:8}
\emph{Trouver la faille.} «~Le teneur de marché a gagné 52.7~\% sur ses fonds propres en 2025, et l'activité de marchés de la banque 16.4~\%~: le market making est plus de
trois fois meilleur métier que la banque.~»
\end{exercise}

\section{Problème~: trois entreprises, une mauvaise année}

\begin{problem}[{Problème du week-end --- trois entreprises, une mauvaise année}]\label{pb:fm:the-economics-of-a-trading-firm:1}
Un conseil d'administration demande comment chacune de trois entreprises cotées survivrait à une année où son \omterm{def:fm:the-economics-of-a-trading-firm:ntr}{revenu net de trading} baisse d'un tiers. On dispose de leurs dépôts.

\medskip\noindent\textbf{Partie I --- Les lignes.}
\begin{enumerate}
\item Définir le \omterm{def:fm:the-economics-of-a-trading-firm:ntr}{revenu net de trading} et les \omterm{def:fm:the-economics-of-a-trading-firm:ntr}{coûts liés au volume}, et expliquer pourquoi le financement fait partie de ces derniers.
\item Calculer le \omterm{def:fm:the-economics-of-a-trading-firm:ntr}{revenu net de trading} du teneur de marché en 2025, et le confronter au chiffre qu'il publie.
\item Donner le résultat opérationnel et la \omterm{def:fm:the-economics-of-a-trading-firm:fixed}{marge opérationnelle} sur \omterm{def:fm:the-economics-of-a-trading-firm:ntr}{revenu net de trading} de chaque entreprise en 2025.
\item Donner le \omterm{def:fm:the-economics-of-a-trading-firm:ntr}{revenu net de trading} et la rémunération par tête du teneur de marché et de la société de gestion.
\item Énoncer ce que la Regulation G et les principes directeurs de l'ESMA exigent d'un \omterm{def:fm:the-economics-of-a-trading-firm:nongaap}{indicateur non GAAP} ou d'une mesure alternative.
\end{enumerate}

\medskip\noindent\textbf{Partie II --- La structure de coûts.}
\begin{enumerate}[resume]
\item Écrire le résultat opérationnel en fonction du \omterm{def:fm:the-economics-of-a-trading-firm:ntr}{revenu net de trading}, avec des coûts fixes $C^{\mathrm{fix}}$ et une part variable $b$.
\item Établir la baisse au seuil de rentabilité et le levier opérationnel, et montrer que l'un est l'inverse de l'autre.
\item Donner la pente estimée $b$ de la rémunération de chaque entreprise avec son erreur type, et indiquer quelles estimations sont fiables.
\item Donner le $C^{\mathrm{fix}}$ de chaque entreprise lorsque la rémunération s'ajuste.
\item Pourquoi le \omterm{def:fm:the-economics-of-a-trading-firm:comp}{ratio de rémunération} de la société de gestion évolue-t-il à l'inverse de son revenu~?
\end{enumerate}

\medskip\noindent\textbf{Partie III --- La mauvaise année.}
\begin{enumerate}[resume]
\item Donner la baisse au seuil de rentabilité de chaque entreprise, avec une rémunération qui s'ajuste puis avec une rémunération maintenue fixe.
\item Donner le levier opérationnel de chaque entreprise dans les deux cas.
\item Quelle part de son résultat opérationnel 2025 chaque entreprise conserve-t-elle après une baisse de 30~\%, avec une rémunération qui s'ajuste~?
\item Quelle entreprise survit à une baisse d'un tiers, et laquelle ne survit qu'en raison de l'ajustement de sa rémunération~?
\item Quelle était la \omterm{def:fm:the-economics-of-a-trading-firm:fixed}{marge opérationnelle} du teneur de marché en 2019, et sa baisse au seuil de rentabilité au revenu de cette année-là~?
\end{enumerate}

\medskip\noindent\textbf{Partie IV --- Le verdict.}
\begin{enumerate}[resume]
\item Donner la pente et la corrélation du \omterm{def:fm:the-economics-of-a-trading-firm:ntr}{revenu net de trading} du teneur de marché et de la société de gestion par rapport à la moyenne annuelle de l'indice VIX.
\item Calculer la \omterm{def:fm:the-economics-of-a-trading-firm:roe}{rentabilité des fonds propres} du teneur de marché en 2025 et en 2023.
\item Pourquoi sa \omterm{def:fm:the-economics-of-a-trading-firm:roe}{rentabilité des fonds propres} et la mesure du segment banque ne désignent-elles pas la même chose~?
\item Énoncer le \emph{résultat nommé}~: la baisse au seuil de rentabilité et le levier opérationnel de chaque entreprise, avec une rémunération qui s'ajuste et avec une rémunération fixe.
\item En deux phrases, que doit faire en premier le conseil d'une entreprise à rémunération rigide et à marge faible~?
\end{enumerate}
\end{problem}

\section{Questions d'entretien}

\begin{interviewq}[$\star$ \iqroles{trader, chercheur}]\label{iq:fm:the-economics-of-a-trading-firm:1}
Qu'est-ce qu'un \omterm{def:fm:the-economics-of-a-trading-firm:comp}{ratio de rémunération}, et pourquoi peut-il augmenter une année où la rémunération baisse~?
\end{interviewq}

\begin{interviewq}[$\star$ \iqroles{trader}]\label{iq:fm:the-economics-of-a-trading-firm:2}
Pourquoi le revenu d'un teneur de marché est-il plus élevé les années volatiles~? Nommez les deux canaux.
\end{interviewq}

\begin{interviewq}[$\star\star$ \iqroles{risque}]\label{iq:fm:the-economics-of-a-trading-firm:3}
Deux entreprises ont le même revenu et les mêmes coûts totaux cette année~; l'une verse 30~\% de chaque dollar marginal en prime, l'autre seulement des salaires. Laquelle
a le levier opérationnel le plus élevé, et laquelle survit à une baisse plus importante~?
\end{interviewq}

\begin{interviewq}[$\star\star$ \iqroles{trader, risque}]\label{iq:fm:the-economics-of-a-trading-firm:4}
Une entreprise a un revenu net de 100, des coûts fixes de 50 et des coûts variables égaux à 20~\% du revenu net. Quel est son résultat opérationnel après une baisse de 40~\% du
revenu~?
\end{interviewq}

\begin{interviewq}[$\star\star$ \iqroles{banque}]\label{iq:fm:the-economics-of-a-trading-firm:5}
Un teneur de marché pour compte propre affiche une \omterm{def:fm:the-economics-of-a-trading-firm:roe}{rentabilité des fonds propres} supérieure à 50~\%~; la division marchés d'une banque affiche 16~\%. Donnez trois raisons pour lesquelles elles
ne peuvent pas être comparées directement.
\end{interviewq}

\begin{interviewq}[$\star\star\star$ \iqroles{chercheur}]\label{iq:fm:the-economics-of-a-trading-firm:6}
Vous régressez la rémunération d'une entreprise sur son revenu pendant six ans et trouvez une pente de $0.29\pm0.09$. Dans quelle mesure lui faites-vous confiance, et qu'est-ce qui d'autre pourrait la produire~?
\end{interviewq}
